\newtoggle{withbonus} \toggletrue{withbonus} % set true to show bonus points in grade table

% setting underlying course name (Grundlagenfach Mathematik, §-Unterricht, \ldots)
\newcommand{\mycourse}{Grundlagenfach Informatik}
% name of the class
\newcommand{\myclass}{Klasse 3g}
% set date
\newcommand{\mydate}{10. Januar 2025}

\title{\textbf{Prüfung: Aus Daten lernen, endliche Automaten}}
\author{Thomas Graf\\
	\mycourse, \myclass}
\date{\mydate}
\maketitle
\thispagestyle{headandfoot}

\begin{center}
	\fbox{\parbox{140mm}{\centering
			\begin{itemize}
				\item Dauer der Prüfung: $45$ Minuten
				      %\item \textbf{Zeige in jeder Aufgabe unbedingt Deine Schritte!}
				\item Erlaubte Hilfsmittel:
				      \begin{itemize}
					      \item Schreibzeug
					      \item Taschenrechner
				      \end{itemize}
			\end{itemize}}}
\end{center}
\vspace{10mm}

\begin{center}
	Vorname: \rule{45mm}{0.8pt}\qquad Nachname: \rule{45mm}{0.8pt}
\end{center}
\vspace{20mm}

\begin{center}
	\iftoggle{withbonus} {
		\combinedgradetable[h][questions]
	} {
		\gradetable[h][questions]
	}
\end{center}
\vspace{15mm}
\begin{center}
	Note:
	\vspace{10mm}

	\rule{8mm}{1.2pt}
\end{center}
%\thispagestyle{empty}
\clearpage

\begin{questions}

	\section*{\textcolor{Green}{Theoretische Informatik}}
	\question[2] Gegeben ist die Sprache $L := \lrc{\lambda, a, ba}$. Bestimmen Sie $L^2$. Wie viele Elemente enthält $L^2$?
	\vspace{5cm}
	\clearpage
	\question{\textbf{Endlichen Automaten zeichnen 1}}

	Gegeben ist die Sprache
	\begin{align*}
		L_1 := \setcm{0x11y}{x,y\in\lrc{0,1}^*}.
	\end{align*}
	\begin{parts}
		\part[1] Nennen Sie ein kürzestes Wort in $L_1$.
		\begin{solution}
			Das Wort $011$ ist \textbf{das} (eindeutige) kürzeste Wort in $L_1$.
		\end{solution}
		\vspace{5cm}
		\begin{solution}
			\begin{figure}[H]
				\centering
				\begin{tikzpicture}[->,>=stealth',shorten >=1pt,node distance=3.2cm,semithick]
					\tikzstyle{every state}=[fill=green!10,draw=black,text=black,minimum size=1.4cm]

					\node[initial,state] (q0) {$q_{0}$};
					\node[state] (q1) [above of=q0] {$q_{1}$};
					\node[state] (q2) [right of=q0] {$q_{2}$};
					\node[state] (q3) [right of=q2] {$q_{3}$};
					\node[state,accepting] (q4) [right of=q3] {$q_{4}$};

					\path (q0) edge [left] node {$1$} (q1)
					(q0) edge [below] node {$0$} (q2)

					(q1) edge [loop above] node {$0,1$} (q1)

					(q2) edge [below] node {$1$} (q3)
					(q2) edge [loop above] node {$0$} (q2)

					(q3) edge [above, bend right] node {$0$} (q2)
					(q3) edge [below] node {$1$} (q4)

					(q4) edge [loop above] node {$0,1$} (q4);
				\end{tikzpicture}
				\caption{Endlicher Automat (in Diagrammform) für die Sprache $L_1$}
				\label{fig:L1}
			\end{figure}
		\end{solution}
		\part[2] Entwerfen Sie einen endlichen Automaten (in Diagrammform), welcher $L_1$ akzeptiert.
	\end{parts}
	\droptotalpoints

	\clearpage \question{\textbf{Endlichen Automaten zeichnen 2}}

	Gegeben ist die Sprache \begin{align*} L_2 := \setcm{x1y}{x,y\in\lrc{0,1}^*, \abs{y} = 1}.
	\end{align*}
	\begin{parts}
		\part[1] Nennen Sie ein kürzestes Wort in $L_2$.
		\begin{solution}
			Die beiden Wörter $10$ und $11$ sind die beiden kürzesten Wörter in $L_2$.
		\end{solution}
		\vspace{5cm}
		\part[3] Entwerfen Sie einen endlichen Automaten (in Diagrammform), welcher $L_2$ akzeptiert.
		\begin{solution}
			\begin{figure}[H]
				\centering
				\begin{tikzpicture}[->,>=stealth',shorten >=1pt,node distance=5.2cm,semithick]
					\tikzstyle{every state}=[fill=green!10,draw=black,text=black,minimum size=1.4cm]

					\node[initial,state] (q0) {$q_{0}$};
					\node[state] (q1) [right of=q0] {$q_{1}$};
					\node[state,accepting] (q2) [right of=q1] {$q_{2}$};
					\node[state,accepting] (q3) [below right of=q1] {$q_{3}$};

					\path (q0) edge [loop below] node {$0$} (q0)
					(q0) edge [above] node {$1$} (q1)

					(q1) edge [below, bend right] node {$0$} (q2)
					(q1) edge [below] node {$1$} (q3)

					(q2) edge [above, bend right] node {$0$} (q0)
					(q2) edge [above] node {$1$} (q1)

					(q3) edge [right] node {$0$} (q2)
					(q3) edge [loop below] node {$1$} (q3);
				\end{tikzpicture}
				\caption{Endlicher Automat (in Diagrammform) für die Sprache $L_2$}
				\label{fig:L2}
			\end{figure}
		\end{solution}
	\end{parts}
	\droptotalpoints
	\clearpage

	\question[2]{\textbf{Anzahl der Elemente in einer Sprache bestimmen}}

	Es sei $\Sigma = \lrc{a, b, c}$ ein Alphabet. Bestimmen Sie die Anzahl der verschiedenen Wörter der Länge $n\in\N$ über $\Sigma$, welche das Teilwort $a$ enthalten.
	\begin{solution}
		Insgesamt gibt es genau $\abs{\Sigma}^n = 3^n$ Wörter über $\Sigma$. Davon enthalten genau $2^n$ Wörter nicht $a$ als Teilwort. Die gesuchte Anzahl ist somit $3^n - 2^n$.
	\end{solution}
	\droptotalpoints
	\clearpage


	\question{\textbf{\red{Nicht reguläre} Sprache}}

	Für ein Wort $a := a_1a_2\ldots a_n$ bezeichnet $a^R := a_na_{n-1}\ldots a_1$ die Umkehrung (\textbf{R}eversal) von
	$a$. Zum Beispiel ist die Umkehrung $w^R$ des binären Worts $w := 11101$ gegeben durch $w^R = 10111$.

	Gegeben sei nun die Sprache
	\begin{align*}
		L_3 := \setcm{wabw^R}{w\in\lrc{a,b}^*}.
	\end{align*}
	\begin{parts}
		\part[1] Nenne genau drei verschiedene Wörter, die in $L_3$ liegen.
		\begin{solution}
			\begin{align*}
				aaabaa, baabab, ab
			\end{align*}
		\end{solution}
		\vspace{3cm}
		\part[3] Zeige, dass $L_3$ \red{nicht regulär} ist.
		\begin{solution}
			Angenommen, $L_3$ sei regulär. Dann existiert ein endlicher Automat
			\begin{align*}
				A = (Q, {a,b}, \delta, q_0, F)
			\end{align*}
			und $A$ akzeptiert $L_3$. Sei jetzt $m := \abs{Q}$, also die Anzahl der verschiedenen Zustände des endlichen Automaten
			$A$. Wir betrachten die $m+1$ verschiedenen Wörter
			\begin{align*}
				a^1ab, a^2ab, \ldots, a^{m+1}ab.
			\end{align*}
			Dies sind mehr Worte (genau ein Wort mehr), als $A$ Zustände hat. Nach dem Schubfachprinzip (pigeonhole principle) existieren
			$i,j\in\lrc{1,2,\ldots,m+1}$ mit $i<j$, sodass
			\begin{align*}
				\hat{\delta}\lr{q_0,a^i} = \hat{\delta}\lr{q_0,a^j}.
			\end{align*}
			Nach der im Unterricht bewiesenen Aussage gilt
			\begin{align*}
				a^iz\in L_3 \iff a^jz\in L_3
			\end{align*}
			für alle $z\in\lrc{a,b}^*$. Doch dies gilt nicht, da zum Beispiel für die Wahl $z:=aba^i$
			\begin{align*}
				a^iz = a^iaba^i\in L_3 \quad\text{aber}\quad a^jz = a^jaba^i\notin L_3
			\end{align*}
			gilt, da $a^i$ zwar die Umkehrung von $a^i$ ist aber nicht von $a^j$. Damit war unsere ursprüngliche Annahme falsch und es existiert kein
			endlicher Automat für die Sprache $L_3$.
		\end{solution}
	\end{parts}
	\droptotalpoints

	\clearpage
	\section*{\textcolor{Green}{Beschreibende Statistik}}
	\question
	\begin{parts}
		\part[1] Nennen Sie ein Beispiel zweier Merkmale aus dem Alltag, die stark positiv korreliert sind.
		\begin{solution}
			Im Unterricht haben wir das Beispiel der starken positiven Korrelation zwischen Körpermasse und Körpergrösse gesehen.
		\end{solution}
		\vspace{4cm}
		\part[1] Nennen Sie ein Beispiel zweier Merkmale aus dem Alltag, die stark negativ korreliert sind.
		\vspace{4cm}
		\begin{solution}
			Zum Beispiel sind Aussentemperatur und Heizkosten stark negativ korreliert.
		\end{solution}
		\part[1] Erläutern Sie den Satz
		\begin{center}
			\textit{Korrelation impliziert nicht Kausalität.}
		\end{center}
		mithilfe eines konkreten Beispiels aus dem Alltag.
		\begin{solution}
			\url{https://www.tylervigen.com/spurious-correlations}
		\end{solution}

		\droptotalpoints
	\end{parts}

	\clearpage
	\question{\textbf{Korrelationskoeffizient bestimmen}}

	Aus der Physik wissen wir, dass die Ausprägungen zweier diskrete Merkmale $x$ und $y$ (die wir exakt messen können) zueinander proportional sind mit Proportionalitätskonstante $\gamma > 0$, also $y=\gamma x$. Sei $n\in\N, n>2$. Wir haben einen Datensatz von $n$ Messpaaren
	\begin{align*}
		(x_1,y_1),(x_2,y_2),...,(x_n,y_n)
	\end{align*}
	von Ausprägungen der proportionalen Merkmale $x$ und $y$ gegeben. Es gilt also $y_i = \gamma x_i$ für alle $i\in\lrc{1,2,...,n}$.
	\begin{parts}
		\part[1] Beweisen Sie die Gleichheit
		\begin{align*}
			\Bar{y} = \gamma\Bar{x}.
		\end{align*}
		\vspace{4cm}

		\begin{solution}
			Wir berechnen
			\begin{align*}
				\Bar{y} = \frac{1}{n}\sum_{i=1}^n y_i = \frac{1}{n}\sum_{i=1}^n \gamma x_i = \gamma\lr{\frac{1}{n}\sum_{i=1}^n x_i} = \gamma\Bar{x}
			\end{align*}
		\end{solution}

		\part[2] Der Korrelationskoeffizient $r$ sei für den gegebenen Datensatz definiert. Bestimmen Sie $r$. Erläutern Sie Ihre Überlegungen.
		\begin{solution}
			Da die Messwerte $x_i$ und $y_i$ proportional zueinander sind mit positiver Proportionalitätskonstante $\gamma>0$, können wir annehmen, dass die beiden Merkmale perfekt positiv korreliert sind mit Korrelationskoeffizient $r = 1$. In der Tat weisen wir diese Vermutung mit der folgenden Berechnung nach. Wir berechnen $r$ für den gegebenen Datensatz. Dabei müssen wir nichts weiter tun, als die gegebenen Grössen in die Definition von $r$ einzusetzen und zu vereinfachen:
			\begin{align*}
				r := & \frac{\displaystyle\sum_{i=1}^n(x_i-\Bar{x})(y_i-\Bar{y})}{\sqrt{\displaystyle\sum_{i=1}^n(x_i-\Bar{x})^2}\cdot\sqrt{\displaystyle\sum_{i=1}^n(y_i-\Bar{y})^2}} =                           \\
				     & \frac{\displaystyle\sum_{i=1}^n(x_i-\Bar{x})(\gamma x_i-\gamma\Bar{x})}{\sqrt{\displaystyle\sum_{i=1}^n(x_i-\Bar{x})^2}\cdot\sqrt{\displaystyle\sum_{i=1}^n(\gamma x_i-\gamma\Bar{x})^2}} = \\
				     & \frac{\displaystyle\sum_{i=1}^n(x_i-\Bar{x})(\gamma(x_i-\Bar{x}))}{\sqrt{\displaystyle\sum_{i=1}^n(x_i-\Bar{x})^2}\cdot\sqrt{\displaystyle\sum_{i=1}^n(\gamma(x_i-\Bar{x}))^2}} =           \\
				     & \frac{\gamma\displaystyle\sum_{i=1}^n(x_i-\Bar{x})^2}{\gamma\displaystyle\sum_{i=1}^n(x_i-\Bar{x})^2} = 1.
			\end{align*}
		\end{solution}
		\droptotalpoints

	\end{parts}

\end{questions}
\clearpage
(leere Seite für Notizen)
\clearpage
(leere Seite für Notizen)

